\documentclass[a4paper,11pt]{article}

\usepackage{amsmath, amssymb}
\usepackage{geometry}
\geometry{margin=20mm}

\usepackage{hyperref}
\usepackage{setspace}
\usepackage{biblatex}
\addbibresource{refs.bib}
\setstretch{0.95}

\title{Create or Curate?\\ How Instructor Framing Shapes the Value of Educational Video\\
\large  1022179 Riki Nishino\\
Advisor:Ian Frank}
\author{}
\date{}

\begin{document}

\maketitle

\section{Research Background}

In recent years, with the promotion of digital transformation (DX) in higher education and the widespread adoption of online learning environments, the use of educational videos has been increasing. Noetel et al. (2021) conducted a systematic review and meta-analysis on the use of video in higher education and demonstrated that educational videos contribute to improved learning outcomes. This suggests that educational videos are widely recognized as effective learning resources in university education.
However, the production of educational videos requires considerable time and technical expertise. Liesa-Orús et al. (2023) reported significant differences in the digital competence of university instructors, indicating that not all faculty members are able to produce their own video materials. Therefore, in educational practice, it is important not only to create new videos but also to select and utilize existing video materials.
Against this background, interest in digital curation in education has been increasing. Deschaine and Sharma (2015) define digital curation as an educational activity involving the selection, organization, and sharing of learning resources that are valuable for learners. In addition, Flintoff et al. (2014) argue that digital curation is an important practice supporting education, research, and professional development.
However, although previous studies have examined the effectiveness of educational videos and the educational value of digital curation, little is known about how learners perceive and evaluate these resources. In particular, even when the same video is used, learners may respond differently depending on whether it is described as “a video created by the instructor” or “a video curated and selected by the instructor.”

\section{Research Purpose}

The purpose of this study is to clarify how instructor-provided information attached to educational videos influences learners’ evaluations of those videos.
Specifically, this study examines whether learners’ perceptions of educational value, usefulness, trustworthiness, satisfaction, and learning motivation differ depending on whether the same instructional video is described as “created by the instructor” or “curated and selected by the instructor.”
Through this, the study aims to investigate the effectiveness of the concept of digital curation in education.

\section{Research Question}

How effective is the concept of “curation” in education?
More specifically, when the same educational video is presented with different descriptions, how does it affect learners’ evaluations of the video?

\section{Related Research}

Research on educational videos has shown that learners’ acceptance and evaluation of videos influence their learning behavior. Küster et al. (2019) demonstrated that perceived usefulness and behavioral intention are important factors in the use of educational videos.
In addition, Shoufan and Mohamed (2022) reviewed the use of YouTube in education and pointed out that, beyond instructors producing their own videos, selecting and integrating existing videos into teaching is also an important educational practice.
Furthermore, Hovland and Weiss (1951) demonstrated that the evaluation of identical information can change depending on perceptions of the source. This finding may also apply to educational videos, suggesting that learners evaluate not only the content of the video but also its perceived origin—whether it is “created by the instructor” or “curated by the instructor”—when judging its value and credibility.
However, few studies have examined the impact of manipulating only the description of the same educational video (i.e., “created” vs. “curated”) on learners’ evaluations.

\section{Research Method}

This study will conduct a small-scale experiment using the same educational video.
Participants will be randomly assigned to one of three conditions:

\begin{itemize}
  \item Creation condition:“The video was created by the instructor for this course.”
  \item Curation condition:“The video was curated by the instructor for this course.”
  \item Neutral condition:No additional information about the origin of the video is provided.
\end{itemize}

To isolate the effect of instructor-provided information, the same video will be used in all conditions.
After viewing the video, participants will complete a comprehension test and a questionnaire. The questionnaire will measure perceived educational value, usefulness, trustworthiness, satisfaction, and learning motivation.
The results will be compared across conditions to analyze how instructor framing influences learners’ perceptions and evaluations of educational videos.

\section{Expected Contributions}

This study does not compare the quality or learning effectiveness of different educational videos themselves, but rather focuses on how instructor framing influences learners’ perceptions.
The findings are expected to contribute to understanding the value of selecting and appropriately contextualizing existing videos in situations where producing new educational videos is difficult.
Furthermore, this study aims to provide evidence that in the use of educational videos, not only “creating” but also “selecting and positioning” resources holds educational significance.


\begin{thebibliography}{99}

\bibitem{noetel2021}
Noetel, M., et al. (2021).
Video improves learning in higher education: A systematic review.
\textit{Review of Educational Research}, 91(2), 204--236.

\bibitem{liesa2023}
Liesa-Orús, M., et al. (2023).
Digital competence of university teachers: A systematic review.
\textit{Education Sciences}, 13(5), 508.

\bibitem{deschaine2015}
Deschaine, M. E., \& Sharma, R. (2015).
The five stages of digital curation: Supporting teaching and learning in online learning environments.
\textit{Journal of Online Learning and Teaching}, 11(3), 398--409.

\bibitem{flintoff2014}
Flintoff, K., Mellow, P., \& Clark, M. (2014).
Digital curation: Opportunities for learning, teaching, research and professional development.
In \textit{Proceedings of ASCILITE 2014}.

\bibitem{kuster2019}
Küster, D., Schumacher, R., \& Ehlers, U.-D. (2019).
The influence of educational videos on learning outcomes in higher education: An empirical study.
\textit{Education and Information Technologies}, 24, 299--320.

\bibitem{shoufan2022}
Shoufan, A., \& Mohamed, F. (2022).
YouTube and Education: A Scoping Review.
\textit{IEEE Access}, 10, 125547--125592.

\bibitem{hovland1951}
Hovland, C. I., \& Weiss, W. (1951).
The influence of source credibility on communication effectiveness.
\textit{Public Opinion Quarterly}, 15(4), 635--650.

\end{thebibliography}

\end{document}